\begin{center}		
\begin{tikzpicture}
		
		
			\begin{axis}[/pgf/number format/.cd, use comma,1000 sep={\,},
				width=13cm,height=13cm,
				xmin=-3.4,xmax=4.4,ymin=-12.7,ymax=5.2,
				xtick={-4,-3,...,5},
				minor x tick num =4,
				ytick={-14,-13,...,5},
				minor y tick num =4,
				tick label style={font=\scriptsize},
				grid=both, axis lines=center,
				grid style={gray!35,line width=0.3pt},
				minor grid style={gray!20,line width=0.2pt},
				enlargelimits=false]
				\addplot[color=blue,domain=-3:4,smooth,samples=200,line width=1.2pt] {-1*(x+1)*(x-3)} node [above right,pos=0.8] {$\mathcal{C}_f$};
				\node[below left] at (axis cs:0,0) {\scriptsize 0};			
				\coordinate[label=135:A] (A) at (axis cs:-2,-5) ;
				\node[draw=black,thick,rotate=45,cross out,inner sep=0pt,minimum width=3pt,minimum height=3pt] at (A) {};
				\coordinate[label=135:C] (C) at (axis cs:-1,1) ;
				\node[draw=black,thick,rotate=45,cross out,inner sep=0pt,minimum width=3pt,minimum height=3pt] at (C) {};
				\coordinate[label=90:B] (B) at (axis cs:1,4) ;
				\node[draw=black,thick,rotate=45,cross out,inner sep=0pt,minimum width=3pt,minimum height=3pt] at (B) {};
				\addplot[color=red,domain=-5:6,smooth,samples=2,line width=1pt] {-5+6*(x+2)} node [above left,pos=0.38] {$(d)$};
				\addplot[color=green4,domain=-5:6,smooth,samples=2,line width=1pt] {4} node [above left,pos=0.65] {$T$};
			\end{axis}
				
		\end{tikzpicture}
\end{center}



On considère la courbe représentative $\mathcal{C}_f$ d'une fonction $f$ définie sur l'intervalle $[-3\,;\,4]$.\par
On note $f'$ la fonction dérivée de la fonction $f$.

On précise que :
\begin{itemize}
\item la droite $(d)$ est tangente à la courbe $\mathcal{C}_f$ au point A et passe par le point C ;
\item la tangente à $\mathcal{C}_f$ au point B est parallèle à l'axe des abscisses.
\end{itemize}

\textbf{Par lecture graphique :}
\begin{enumerate}
\item Déterminer $f(-2)$ et $f(1)$.
\item Déterminer $f'(-2)$ et $f'(1)$.
\item Résoudre l'équation $f(x)=0$.
\item Dresser le tableau de variation de $f$.
\end{enumerate}

\textbf{On admet que $f(x)=-x^2+2x+3$ sur l'intervalle $[-3\,;\,4]$.}
\begin{enumerate}[resume]
\item Retrouver par le calcul les résultats de la question 1.
\item
\begin{enumerate}[label=\alph*)]
\item Déterminer $f'(x)$ puis retrouver les résultats de la question 2.
\item Vérifier par le calcul que $f(x)=(x+1)(-x+3)$.\par Retrouver les résultats de la question 3.
\item Dresser le tableau de signes de $f'(x)$ sur l'intervalle $[-3\,;\,4]$.\par En déduire les variations de la fonction $f$ sur cet intervalle.
\end{enumerate}
\end{enumerate}

\bigskip
